\section{Option Cashflow and Risk Framework}\label{sec:cashflow}

The central modeling step is to express every option as a NAV-normalized cashflow. Let
\(W_t\) be beginning portfolio NAV, \(q_{i,t}\) the signed premium weight assigned to option
\(i\in\{1,\dots,n\}\), and \(C_{i,t}>0\) the decision-date mark. The contract count is
\[
  N_{i,t}=\frac{q_{i,t}W_t}{C_{i,t}}.
\]
Ignoring transaction costs and slippage for the gross theory object, but not ignoring the
option purchase price, the next-period NAV contribution is
\[
  \frac{\Delta W_{i,t+1}}{W_t}
  =q_{i,t}\frac{C_{i,t+1}-C_{i,t}}{C_{i,t}}
  =q_{i,t}\,r_{i,t+1},
  \qquad
  r_{i,t+1}:=\frac{C_{i,t+1}-C_{i,t}}{C_{i,t}},
\]
which defines the one-period option mark return \(r_{i,t+1}\). These are premium weights. A
long option whose mark falls to zero loses exactly the paid premium weight. A short option
is scaled by the premium it sold and then pays the future mark or payoff liability. Summing
over contracts and the residual cash account turns the contribution into an exact NAV
recursion.

\begin{proposition}[NAV accounting identity]\label{prop:nav}
Let \(q_t\in\R^n\) be the premium-weight vector, \(r_{t+1}\) the option mark-return vector,
and \(r_{f,t}\) the one-period financing rate earned on the residual cash
\(W_t(1-\one^\top q_t)\). Then
\[
  \frac{W_{t+1}}{W_t}-1
  =q_t^\top r_{t+1}+\bigl(1-\one^\top q_t\bigr)r_{f,t},
  \qquad
  \frac{W_{t+1}}{W_t}-1-r_{f,t}
  =q_t^\top\bigl(r_{t+1}-r_{f,t}\one\bigr).
\]
\end{proposition}

\begin{proof}
At \(t+1\) the option book is worth
\(\sum_i N_{i,t}C_{i,t+1}=W_t\sum_i q_{i,t}(1+r_{i,t+1})\) and the cash account is worth
\(W_t(1-\one^\top q_t)(1+r_{f,t})\). Adding the two, dividing by \(W_t\), and subtracting
\(1\) (and then \(r_{f,t}\)) gives both displays.
\end{proof}

\begin{remark}[Zero-financing numeraire]\label{rem:zero_financing}
The baseline empirical run sets \(r_{f,t}=0\): residual cash and intra-month expiry
proceeds earn nothing. Every Sharpe object reported in the paper is therefore a raw-return
Sharpe ratio measured against a zero threshold, not an excess-of-cash Sharpe ratio. With
\(r_{f,t}\neq0\), Proposition~\ref{prop:nav} shows the correct objective is the excess
return \(q^\top(r_{t+1}-r_{f,t}\one)\), and the unconstrained maximum-Sharpe direction
(Proposition~\ref{prop:maxsharpe}) becomes
\(\Sigma_{O,t}^{-1}(\mu_{O,t}-r_{f,t}\one)\). Over the 2021--2026 sample, short rates were
materially positive for much of the window, so the zero-financing convention forgoes
collateral yield in the numerator while also applying a zero rather than a positive hurdle
in the excess return; the two effects act in opposite directions. All strategies and
benchmarks are reported under the same convention, so comparisons remain internally
consistent, and readers can restate any row against a cash hurdle using the identity above.
\end{remark}

Two notational conventions apply throughout. The operator \(\Delta\) applied to a state
variable always denotes a one-rebalance-period difference, for example
\(\Delta S_u=S_{u,t+1}-S_{u,t}\) and \(\Delta\sigma_i=\sigma_{i,t+1}-\sigma_{i,t}\); the
Black--Scholes and Black--76 hedge ratios carry a superscript,
\(\Delta^{\mathrm{BS}}_i\) for equity options and \(\Delta^{VX}_j\) for VIX options, so the
two uses never collide. The underlying simple return is \(R_u=\Delta S_u/S_u\).

For a call or put on equity underlying \(u\), the mark change is decomposed by the same
Ito--Taylor accounting that underlies Black--Scholes--Merton option valuation
\citep{blackscholes1973,merton1973,shreve2004stochasticii}:
\[
  \Delta C_i
  =
  \Delta^{\mathrm{BS}}_i\,\Delta S_u
  +\frac12\Gamma_i(\Delta S_u)^2
  +\nu_i\,\Delta\sigma_i
  +\Theta_i\,\Delta t
  +\eta_i,
\]
where \(\eta_i\) is \emph{defined} as the exact remainder, so the display holds with
equality rather than approximately. The paper does not require Black--Scholes to be a
globally correct pricing model for American single-name options. It uses the weaker local
statement: an option mark can be expanded in spot, volatility, and time, with vanna, volga,
higher-order terms, early-exercise premium movement, and marking error left in \(\eta_i\).

Dividing by the decision-date mark \(C_{i,t}\) gives the return exposure map
\[
 r_{i,t+1}
 =
 \underbrace{\frac{\Delta^{\mathrm{BS}}_i S_u}{C_{i,t}}}_{b^\Delta_{i,t}}R_u
 +\underbrace{\frac{\Gamma_i S_u^2}{2C_{i,t}}}_{b^\Gamma_{i,t}}R_u^2
 +\underbrace{\frac{\nu_i}{C_{i,t}}}_{b^\nu_{i,t}}\Delta\sigma_i
 +\underbrace{\frac{\Theta_i\Delta t}{C_{i,t}}}_{\theta^{\ast}_{i,t}}
 +\epsilon_{i,t+1},
 \qquad
 \epsilon_{i,t+1}=\frac{\eta_i}{C_{i,t}},
\]
again with equality because the scaled remainder is retained. The contract-level implied-vol
change is mapped to its underlying's at-the-money factor through
\(\Delta\sigma_i=\beta^{\sigma}_{i,t}\,\Delta\sigma^{\mathrm{ATM}}_{u(i)}
+(\text{smile remainder})\), with the smile remainder absorbed into \(\epsilon_{i,t+1}\).
The coefficients are dimensionless option-return exposures. Deep out-of-the-money options
can have large return exposure because the premium denominator is small. That is not a
numerical nuisance. It is the reason option-only portfolios need explicit gross,
per-contract, Greek, beta, and stress budgets.

\begin{remark}[Unit conventions]\label{rem:units}
Marks \(C_{i,t}\) and Greeks are quoted per share; contract-level quantities carry the
100-share multiplier, which cancels in every premium-weight ratio above. Volatilities are
in decimal units (one vol point \(=0.01\)), \(\Theta_i\) is annualized, and \(\Delta t\) is
the rebalance horizon in year fractions. Contract counts \(N_{i,t}=q_{i,t}W_t/C_{i,t}\) are
fractional in the theory object; rounding to integer contracts perturbs each premium weight
by at most \(C_{i,t}/W_t\) per contract, which the implementation treats as part of the
cost and capacity layer rather than of the risk model.
\end{remark}

Stacking contracts gives the conditional factor form. The factor vector collects underlying
returns, centered squared returns, at-the-money implied-volatility shocks, skew and
tail-slope shocks, and the matched VX-forward returns with their centered squares,
\[
\begin{aligned}
  f_{t+1}
  =\Bigl(\,&R_{u},\;
  R_{u}^{2}-\E_t\bigl[R_{u}^{2}\bigr],\;
  \Delta\sigma^{\mathrm{ATM}}_{u},\;
  \Delta\mathrm{skew}_{u},\;
  \Delta\mathrm{tail}_{u},\\
  &R^{VX}_{T},\;
  \bigl(R^{VX}_{T}\bigr)^{2}-\E_t\bigl[(R^{VX}_{T})^{2}\bigr]
  \Bigr)_{u\in\mathcal U,\;T\in\mathcal T},
\end{aligned}
\]
where \(\mathcal U\) is the set of equity underlyings and \(\mathcal T\) the set of
VX-forward maturities. The matching rows of \(B_t\) are
\(\bigl(b^\Delta_{i,t},\,b^\Gamma_{i,t},\,b^\nu_{i,t},\,b^{skew}_{i,t},\,b^{tail}_{i,t},\,
b^{\Delta,VX}_{i,t},\,b^{\Gamma,VX}_{i,t}\bigr)\). Writing
\(\Omega_t=\Var_t(f_{t+1})\) for the conditional factor covariance,
\[
  r_{O,t+1}=\mu_{O,t}+B_t f_{t+1}+\varepsilon_{t+1},
  \qquad
  \Sigma_{O,t}=B_t\Omega_tB_t^\top+\Sigma_{\varepsilon,t},
\]
where \(\Sigma_{\varepsilon,t}=\Var_t(\varepsilon_{t+1})\) is the residual option-return
covariance. The residual term is part of the model. Omitting it would assume that Greeks
span all marking risk. The covariance identity requires one formal statement about what
\(\varepsilon_{t+1}\) is.

\begin{assumption}[Projection-residual definition of the residual]\label{assump:projection}
\(\Omega_t=\Var_t(f_{t+1})\succ0\), and \(B_t\) and \(\varepsilon_{t+1}\) are defined by
the population linear projection of the centered option return on the factors,
\[
  B_t:=\Cov_t\bigl(r_{O,t+1},f_{t+1}\bigr)\,\Omega_t^{-1},
  \qquad
  \varepsilon_{t+1}:=r_{O,t+1}-\mu_{O,t}-B_t f_{t+1}.
\]
Consequently \(\Cov_t(f_{t+1},\varepsilon_{t+1})=0\) holds by construction, and
\(\Sigma_{O,t}=B_t\Omega_tB_t^\top+\Sigma_{\varepsilon,t}\) is an identity
(Proposition~\ref{prop:cov_psd}), not an approximation.
\end{assumption}

\begin{remark}[The Greek map is an estimator of \(B_t\), not its definition]\label{rem:greekmap}
The analytic loadings in the return exposure map are the structural estimator of the
projection coefficient \(B_t\). The wedge between the two---vanna and volga terms, the
higher-order Taylor remainder, and early-exercise premium movement in American
contracts---is measured empirically in the approximation-diagnostics table
(Table~\ref{tab:approx_app}). If \(B_t\) were instead \emph{fixed} at the analytic Greek
loadings, the residual would in general correlate with the factors and the covariance
would acquire the omitted cross term
\(B_t\Sigma_{f\varepsilon,t}+\Sigma_{f\varepsilon,t}^\top B_t^\top\). The paper therefore
treats the Greek map as an estimator whose adequacy is reported, rather than assuming the
cross term away.
\end{remark}

\begin{remark}[Rebalance-horizon approximation]\label{rem:horizon}
The Taylor decomposition is applied at the monthly rebalance horizon with decision-date
Greeks held fixed; it is a working risk model whose adequacy is tested empirically in the
approximation diagnostics, not an asymptotic small-\(\Delta t\) statement.
\end{remark}

VIX options require a separate convention. They are not options on an equity spot. The
traded forward-risk anchor is the matched VX futures curve, and Black's futures-option
logic is the natural local model \citep{black1976commodity}. For VIX option \(j\) with
matched forward \(F^{VX}_{T_j}\), the local decomposition is
\[
  \Delta C^{VIX}_j
  =
  \Delta^{VX}_j\,\Delta F^{VX}_{T_j}
  +\frac12\Gamma^{VX}_j\bigl(\Delta F^{VX}_{T_j}\bigr)^2
  +\nu^{VIX}_j\,\Delta\sigma^{VIX}_j
  +\Theta^{VIX}_j\,\Delta t
  +\eta^{VIX}_j,
\]
with \(\eta^{VIX}_j\) the exact remainder. Dividing by the mark \(C_{j,t}\) gives the
premium-normalized form that enters the factor model,
\[
  r_{j,t+1}
  =b^{\Delta,VX}_{j,t}R^{VX}_{T_j}
  +b^{\Gamma,VX}_{j,t}\bigl(R^{VX}_{T_j}\bigr)^2
  +b^{\nu,VX}_{j,t}\,\Delta\sigma^{VIX}_j
  +\theta^{\ast}_{j,t}
  +\epsilon_{j,t+1},
\]
\[
  b^{\Delta,VX}_{j,t}=\frac{\Delta^{VX}_j F^{VX}_{T_j}}{C_{j,t}},\qquad
  b^{\Gamma,VX}_{j,t}=\frac{\Gamma^{VX}_j\bigl(F^{VX}_{T_j}\bigr)^2}{2C_{j,t}},\qquad
  b^{\nu,VX}_{j,t}=\frac{\nu^{VIX}_j}{C_{j,t}},
\]
where \(R^{VX}_{T_j}=\Delta F^{VX}_{T_j}/F^{VX}_{T_j}\) is the matched VX-forward return
entering \(f_{t+1}\) on equal footing with the equity underlyings, and
\(\epsilon_{j,t+1}=\eta^{VIX}_j/C_{j,t}\). VIX, VVIX, and term-structure indices are useful
state variables; VX futures and VIX options are the traded volatility instruments. Exact
VRO/SOQ settlement is required before VIX option expiry P\&L can be promoted from
settlement proxy to headline listed-settlement evidence.
